\color{unimain}{\subsection{Настройка каналов напряжения}\label{func:nkn}}
\color{black}

\begin{enumerate}
\item
Для подключения вторичных цепей напряжения от измерительных трансформаторов напряжения в устройстве предусмотрены следующие аналоговые входы:
\begin{itemize}
	\item 3 входа для подключения цепей напряжения от основной вторичной обмотки ТН (Ua, Ub, Uc);
	\item 3 входа для подключения цепей напряжения от дополнительной вторичной обмотки ТН (Uни[Uнф], Uик[Uфк], Uиф[Uнк]);
	\item 1 вход для подключения цепи отбора напряжения ЛЭП (Uотб ЛЭП);
	\item 1 вход для подключения цепи отбора напряжения Э1 (Uотб Э1);
	\item 1 вход для подключения цепи отбора напряжения Э2 (Uотб Э2).
\end{itemize}

\item
Рабочий ТН, трехфазная система которого используется для работы защит и автоматики может быть установлен как на линии, так и на шинах. Для корректной работы логики переключатель \textbf{«Место ТН»} должен быть установлен в соответствующее положение.

\item
Аналоговые входы отбора напряжения ЛЭП, Э1 и Э2, на которые могут быть поданы напряжения от однофазного ТН или ШОН, предназначены для целей АПВ и АУ. 

\item
Схема настройки аналоговых значений $U_{\textsl{ЛЭП}}$, $U_{\textsl{Э1}}$, $U_{\textsl{Э2}}$ в зависимости от схемы РУ приведена на рисунке \ref{pict:nkn}.

\medskip

\begin{figure}[H]
\centering
\includegraphics[width=0.8\textwidth,height=0.8\textheight,keepaspectratio]{\fbpath/03.01 ЛВ/22. НКН/img/nkn.pdf}
\captionof{figure}{Функциональная схема настройки аналоговых значений $U_{\textsl{ЛЭП}}$, $U_{\textsl{Э1}}$, $U_{\textsl{Э2}}$}\label{pict:nkn}
\end{figure}

\medskip

\item
Если, например, рабочий ТН установлен на шинах, то значению $U_{\textsl{Э1}}$ присваивается одно из фазных либо линейных напряжений, выбираемое переключателем \textbf{«Фаза подключения Uэ1»}. Переключатель \textbf{«Напр. Uэ1»} в таком случае должен быть в положении \textit{«Фаза подкл.»}. А значению $U_{\textsl{ЛЭП}}$ присваивается напряжение, поступающее с аналогового входа отбора напряжения ЛЭП (Uотб ЛЭП), скорректированное по амплитуде и фазе соответствующими уставками \textbf{«К подстройки Uотб ЛЭП»} и \textbf{«Ф подстройки Uотб ЛЭП»} по следующим формулам:

\begin{align}
\vec{U}_{\text{отб лэп корр}} = \vec{U}_{\text{отб лэп}} \cdot K_{\text{лэп}} \cdot e^{j\phi_{\text{лэп}}}
\end{align}
\begin{eqexpl}[7mm]
  \item{$ K_{\text{лэп}} $} коэффициент подстройки Uотб ЛЭП по амплитуде;
	\vspace{2pt}
  \item{$ \phi_{\text{лэп}} $} угол подстройки Uотб ЛЭП по фазе.
\end{eqexpl}

\item
Если рабочий ТН установлен на ЛЭП, то наоборот: фаза подключения выбирается для $U_{\textsl{ЛЭП}}$, а значению $U_{\textsl{Э1}}$ присваивается скорректированное напряжение с аналогового входа отбора напряжения Э1 (Uотб Э1). 

\end{enumerate}